\documentclass[a0,portrait]{a0poster}
\usepackage{MLZposter}
\usepackage{caption}
\usepackage{showyourwork}
\usepackage[utf8]{inputenc}
% \usepackage{html}
% \stylesheet{style.css}
% \usepackage[colorlinks=true, citecolor=blue, linkcolor=blue, urlcolor=blue]{hyperref}
% \hypersetup{
%     colorlinks=true,
%     urlcolor=blue,
%     citecolor=blue,
%     linkcolor=blue
% }
% \usepackage[backend=biber,style=numeric,sorting=none]{biblatex}
% \bibliographystyle{unsrt}
% \bibliography{bib}
\usepackage[numbers, sort&compress]{natbib} 


%%%%%%%%%%%%%%%%%%%%%%%%%%%%%%%%%%%%%%%%%%%%%%%%%%%%%%%%%%%%%%%%%%%%%%%%%%%%%%%%
\begin{document}

%%% Top box of poster showing the Logos of MLZ, ... %%%
\begin{tikzpicture}[remember picture,overlay]
  \draw[color=mlzblue,line width=6pt,rounded corners=20pt] (current page.north west) ++(4.5,0) -- +(0,-6.9) -- +(84,-6.9);
  \node[color=mlztextgray,text width=12cm] at (72.8,-1){
\includegraphics[height=5.0cm]{figures/mlz_blue_rgb.png}};
  \node[color=mlztextgray,text width=9cm] at (9.5,-1){
\includegraphics[height=5.0cm]{figures/hereon.png}};     
\end{tikzpicture}

%%% Title block %%%
\begin{textblock}{86}(7.5,9)
   \begin{center}
   \color{mlzblue}
    {\bfseries\veryHuge Workflow for creating a reproducible research article using showyourwork \\[1.2ex]}
    {\LARGE \underline{Arnab Majumdar}\textsuperscript{1,2}, Sebastian Busch\textsuperscript{1,2}}\\[1.5ex]
    {\large\textsuperscript{1}German Engineering Materials Science Center (GEMS) at Heinz Maier-Leibnitz Zentrum (MLZ), Helmholtz-Zentrum Hereon,\\ \textsuperscript{2}Institute of Materials Physics, Helmholtz-Zentrum Hereon.}\\[.9ex]
   \end{center}
\end{textblock}


%%%%%%%%%%%%%%%%%%%%%%%%%%%%%%%%%%%%%%%%%%%%%%%%%%%%%%%%%%%%%%%%%%%%%%%%%%%%%%%%
\begin{textblock}{86}(7.5,19.5)
\SECTION{Aim of daphne}
\begin{tikzpicture}
	\anode{0}{0}{74cm}{Ensuring reproducibility in scientific publications remains a major challenge for photon and neutron experiments \cite{peng2021reproducible}. To address this, DAPHNE aims to develop a bidirectional workflow that links raw data and publications via processing scripts used for data reduction and analysis scripts used for data analysis, as shown below.
    
    \begin{figure}
    {
        \begin{center}
        \includegraphics[width=0.7\textwidth]{figures/daphne_proposal.png}
        \label{fig:daphne_proposal}   
        \caption{Snapshot from daphne proposal~\cite{barty2023daphne}: [Top] Current standard of publication without any raw data, [middle] currently possible situation of linking publication to raw data from facility, and [bottom] inclusion of processing (or reduction) and analysis steps between raw data and publication with bidirectional links.}
        \end{center}
    }
    \end{figure}
    }
\end{tikzpicture}

\end{textblock}

%%%%%%%%%%%%%%%%%%%%%%%%%%%%%%%%%%%%%%%%%%%%%%%%%%%%%%%%%%%%%%%%%%%%%%%%%%%%%%%%
\begin{textblock}{86}(7.5,43)
\SECTION{Showyourwork philosophy for reproducing scientific article}
\begin{tikzpicture}
	\anode{0}{0}{74cm}{
    \emph{Showyourwork!}~\cite{Luger2021} is an open-source reproducibility framework designed to enable transparent and fully reproducible scientific research. It organizes analysis scripts, software dependencies, and the manuscript within a GitHub repository, allowing any independent researcher to clone the project and regenerate the complete article. The manuscript preparation can be integrated with Overleaf, enabling collaborative writing and seamless synchronization between analysis and documentation.

    Once reproduced, figures can be modified directly by editing the underlying scripts that generate them. Raw datasets are typically accessed from a Zenodo repository with read-only permissions, while intermediate and processed data products can be automatically uploaded to a separate Zenodo repository with read–write access.

    \begin{figure}
    {
        \script{overview.py}
        \begin{center}
        \includegraphics[width=40cm, trim=0cm 1.3cm 0cm 1cm, clip]{figures/overview.pdf}
        \caption{Framework of \emph{showyourwork!}: GitHub hosts analysis scripts, environment files, dependency rules, and the manuscript; Zenodo stores datasets; and Overleaf enables manuscript editing. Commits trigger automated figure generation and manuscript compilation via GitHub Actions, producing a final version and an arXiv-ready tarball. \label{fig:fig_flw_chart}}    
        \end{center}
    }
    \end{figure}
	}
\end{tikzpicture}
\end{textblock}


%%%%%%%%%%%%%%%%%%%%%%%%%%%%%%%%%%%%%%%%%%%%%%%%%%%%%%%%%%%%%%%%%%%%%%%%%%%%%%%%
\begin{textblock}{86}(7.5,64.5)
\SECTION{Demonstration of reproducible workflow using showyourwork}
\begin{tikzpicture}
    \anode{0}{0}{74cm}{
        This poster was created using the \emph{showyourwork!} software and is available under Creative Commons Attribution 4.0 International license. The repository is available at \href{https://github.com/majumdararn/syw-poster-dam}. The figure \ref{fig:daphne_proposal} 
        
        The copy was created using the reproducibility software showyourwork!, which allows continuous integration to programmatically download the data from
zenodo.org, create the figures, and compile the manuscript.
Each figure caption in the electronic copy contains two links:
one to the dataset stored on Zenodo used in the corresponding figure, and the other to the script used to make the figure. The release 1.0.0 of the git repository allows anyone to
rebuild the entire manuscript. The datasets are available at
https://doi.org/10.5281/zenodo.16994688, under the
Creative Commons Attribution 4.0 International license.
    }
	\anode{0}{5}{24cm}{
    \begin{figure}
    {
        \script{raw_data_plot.py}
        \begin{center}
        \includegraphics[width=10cm, trim=0cm 1.3cm 0cm 1cm, clip]{figures/raw_data.pdf}
        \caption{Raw data: 1D sinusoinal curve with rotation\label{fig:fig_raw_data}}    
        \end{center}
    }
    \end{figure}
	}
    \anode{24}{5}{24cm}{
    \begin{figure}
    {
        \script{red_data_plot.py}
        \begin{center}
        \includegraphics[width=10cm, trim=0cm 1.3cm 0cm 1cm, clip]{figures/red_data.pdf}
        \caption{Reduce data: 1D sinusoinal curve without rotation\label{fig:fig_red_data}}    
        \end{center}
    }
    \end{figure}
    }
    \anode{48}{5}{24cm}{
    \begin{figure}
    {
        \script{ana_data_plot.py}
        \begin{center}
        \includegraphics[width=10cm, trim=0cm 1.3cm 0cm 1cm, clip]{figures/ana_data.pdf}
        \caption{data analysis: angular freq and phase shift (expected vs fitted)\label{fig:fig_ana_data}}    
        \end{center}
    }
    \end{figure}
	}
\end{tikzpicture}
\end{textblock}


%%%%%%%%%%%%%%%%%%%%%%%%%%%%%%%%%%%%%%%%%%%%%%%%%%%%%%%%%%%%%%%%%%%%%%%%%%%%%%%%
\begin{textblock}{86}(7.5,85)
\SECTION{References}
\begin{tikzpicture}
	\anode{0}{0}{74cm}{
    \renewcommand{\refname}{}   % remove heading
    \vspace{-1.5em}                    % remove extra vertical space (adjust as needed)
    \bibliographystyle{unsrtnat}
    \bibliography{bib}
	}
\end{tikzpicture}
\end{textblock}

%%%%%%%%%%%%%%%%%%%%%%%%%%%%%%%%%%%%%%%%%%%%%%%%%%%%%%%%%%%%%%%%%%%%%%%%%%%%%%%%

\begin{tikzpicture}[remember picture,overlay]
  \draw[color=mlzblue,line width=6pt,rounded corners=20pt] (current page.south west) ++(0,5.2) -- +(77.5,0) -- +(77.5, -6);
\end{tikzpicture}

\begin{textblock}{86}(7.5,93.9)
\begin{tikzpicture}
  \anode{0}{0}{5cm}{\textbf{Contact:}}
  \anode{4}{0.0}{20cm}{Name Arnab Majumdar\\ E-mail: arnab.majumdar@hereon.de\\ phone: +49.89.158860.771};
  \anode{25}{0.0}{20cm}{Name Sebastian Busch\\ E-mail: sebastian.busch@hereon.de\\ phone: +49.89.158860.764};
  \anode{66}{0}{5cm}{
\includegraphics[height=3.5cm]{figures/gems.png}};
\end{tikzpicture}
\end{textblock}

\end{document}

